% 实验章节

\section{实验}
\label{sec:experiments}

\subsection{实验设置}

所有实验在一台配备Intel Core i7-13700 CPU和NVIDIA GeForce RTX 4060 (8GB) GPU的计算机上进行，深度学习框架为PyTorch 2.11.0+cu128。为保证可复现性，所有随机种子统一设置为42，并在每个模型训练前显式重置所有随机状态（random、numpy、torch、cuda）。FT-Transformer和DeepFM因rtdl和deepctr-torch在当前环境下不适配，使用手写实现。

\subsubsection{XGBoost基线配置}

XGBoost：300轮提升，最大深度8，学习率0.05，scale\_pos\_weight根据训练集正负比自动计算。为充分利用数据，XGBoost在训练集和验证集的合并数据上训练；深度模型仅使用训练集训练并通过验证集早停。

\subsection{模型性能对比}

表\ref{tab:model_comparison}汇总了全部六个模型（三个基座+三个蒸馏变体）在七项评估指标上的定量结果，按教师、基线、蒸馏变体分组排列。所有指标均基于各自最优分类阈值（通过验证集F1最大化确定）计算得出。

\begin{table}[htbp]
\centering
\small
\setlength{\tabcolsep}{4pt}
\caption{全部模型性能对比与蒸馏定量结果。按类别分组：教师$\to$基线$\to$蒸馏变体。FTT+KD（3A only）为消融对照，隔离TreeReg的效应。加粗为各指标最优值。}
\label{tab:model_comparison}
\begin{tabular}{lccccccc}
\toprule
\textbf{模型} & \textbf{AUPRC} & \textbf{ROC-AUC} & \textbf{F1} & \textbf{Prec} & \textbf{Rec} & \textbf{Cost} & \textbf{ECE} \\
\midrule
XGBoost（教师）   & 0.9331 & 0.9979 & 0.8750 & 0.9479 & 0.8125 & 215 & 0.0006 \\
DeepFM（基线）    & 0.9336 & 0.9979 & 0.8938 & 0.8860 & 0.9018 & 123 & 0.0394 \\
FT-Transformer（基线） & 0.9179 & 0.9925 & 0.9196 & 0.9196 & 0.9196 & 99 & 0.0120 \\
\addlinespace
FTT + KD（3A only）& 0.9347 & 0.9953 & \textbf{0.9279} & 0.9364 & 0.9196 & \textbf{97} & 0.0007 \\
FTT + KD + TreeReg（3A+3B）& 0.8974 & 0.9945 & 0.8848 & 0.9143 & 0.8571 & 169 & \textbf{0.0010} \\
DeepFM + KD（3A only）& \textbf{0.9495} & \textbf{0.9988} & 0.9058 & 0.9099 & 0.9018 & 120 & 0.0008 \\
\bottomrule
\end{tabular}
\end{table}

如表\ref{tab:model_comparison}所示，模型呈现出清晰的三层性能梯度：

\textbf{（1）蒸馏模型占据性能前二。}DeepFM+KD以AUPRC=0.9495位列榜首（超越教师XGBoost +1.64pp），FTT+KD以0.9347紧随其后（+1.68pp vs. FTT基线）。两者均取得近乎完美的校准（ECE$\le$0.001）和较低的业务成本（分别为120和97）。这一结果推翻了"深度模型在表格数据上无法超越树模型"的刻板印象——通过知识蒸馏注入教师的决策边界知识，深度模型不仅弥合了归纳偏置差距，甚至实现了反超。

\textbf{（2）基座模型中的异构性能格局。}DeepFM（AUPRC=0.9336）与XGBoost（0.9331）差距仅0.0005——在单一随机种子下统计不可区分——验证了显式二阶交互在PCA空间中具有与轴对齐分裂同等有效的归纳偏置\cite{grinsztajn2022tree}。FT-Transformer经Sparsemax替换Softmax后AUPRC达0.9179（原Softmax版本为0.8738，+4.4pp），ECE从0.2328降至0.0120（$\times$20改善）。但DeepFM的内积交互$\langle v_i, v_j \rangle$直接建模特征对的协同效应，不依赖"哪些特征更重要"的全局判断，在PCA匿名化空间中仍比FT-Transformer的全局自注意力更为稳健。

\textbf{（3）TreeReg的负向效应。}FTT+KD+TreeReg以AUPRC=0.8974垫底——不仅低于FTT+KD（-3.73pp），甚至低于未蒸馏的FTT基线（-2.05pp）。这一结果将因果链条锁定于TreeReg（方案3B）：特征重要性正则化强制嵌入层模仿XGBoost的轴对齐特征偏好，与Sparsemax注意力机制对自由特征交互学习的需求形成结构性拮抗。需指出，FTT+KD+TreeReg的ECE=0.0010仍为优秀（软标签的校准平滑效应在机制上与TreeReg解耦），但其AUPRC损失远超KD收益。

\textbf{（4）ROC-AUC区分度不足，AUPRC至关重要。}全部模型的ROC-AUC取值区间极窄（0.9925--0.9988，极差0.0063），而AUPRC区间达0.0521（0.8974--0.9495），区分度约8.3倍。这一差距在不平衡数据上（欺诈率约0.5\%）符合预期：ROC-AUC以真负率为横轴，被海量正常样本主导，而AUPRC直接衡量模型在正类上的排序能力\cite{dastidar2024survey}。FTT+KD+TreeReg的AUPRC（0.8974）与DeepFM+KD（0.9495）之间0.0521的实际性能差距在ROC-AUC上仅表现为0.0043——若仅报告ROC-AUC，近半的性能差异将被掩盖。

\textbf{（5）成本-性能非单调。}FTT+KD以总成本97取得全局最低（FTT基线为99），而AUPRC最高的DeepFM+KD成本为120，XGBoost为215（全局最高）。AUPRC最优$\neq$成本最优，模型选择应嵌入业务成本矩阵综合权衡。

图\ref{fig:pr_curves}和图\ref{fig:roc_curves}分别展示了全部六个模型的PR曲线和ROC曲线。

\begin{figure}[htbp]
    \centering
    \includegraphics[width=0.78\textwidth]{pr_curves.pdf}
    \caption{全部六个模型的Precision-Recall曲线，图例标注AUPRC值。DeepFM+KD和FTT+KD在召回区间0.6--0.9保持精确率优势，FTT+KD+TreeReg的高召回端急剧下降反映了TreeReg对模型的破坏性影响。}
    \label{fig:pr_curves}
\end{figure}

图\ref{fig:pr_curves}揭示了两层信息。横向上，六个模型的PR曲线构成了清晰的三组群：（a）DeepFM+KD和FTT+KD位于最上方，其在0.4--0.9召回区间内均保持最高的精确率；（b）DeepFM和XGBoost在中段召回区域几乎重合（差距$<$0.005），FT-Transformer略低但差距有限；（c）FTT+KD+TreeReg在所有召回端均低于其他模型，尤其在高召回端（$>$0.85）精确率急剧下降。纵向上，所有曲线在高召回区域（$>$0.95）均急剧下降，反映了极度不平衡数据下精确率与召回率的内在张力：检出几乎所有的欺诈交易必然伴随大量误报。

\begin{figure}[htbp]
    \centering
    \includegraphics[width=0.78\textwidth]{roc_curves.pdf}
    \caption{全部六个模型的ROC曲线，图例标注AUC值，黑色虚线为随机分类器基准。六条曲线视觉差异极小，进一步验证了ROC-AUC在不平衡分类问题上的固有局限。}
    \label{fig:roc_curves}
\end{figure}

与PR曲线形成鲜明对比，图\ref{fig:roc_curves}中六条ROC曲线几乎完全重叠，肉眼无法区分模型优劣——即便是AUPRC差距达0.0521的DeepFM+KD与FTT+KD+TreeReg，在ROC图上也无明显分离。这直观展示了ROC-AUC在不平衡分类问题上的固有缺陷：ROC以真负率（TNR）为横轴，当正类样本仅占约0.5\%时，即便模型将所有样本预测为正常，TNR仍接近1.0，导致AUC虚高。全部模型的ROC-AUC极差仅0.0063（0.9925--0.9988），而AUPRC极差达0.0521——区分度差距约8.3倍。这一对比为本文以AUPRC为主要评估指标的实验设计提供了充分的实证支持。

\paragraph{主要发现} PR曲线与ROC曲线的对比实验表明：（1）蒸馏模型（DeepFM+KD、FTT+KD）在宽召回区间内全面超越基座模型，FTT+KD+TreeReg垫底；（2）六条ROC曲线几乎重叠，AUPRC极差（0.0521）约8.3倍于ROC-AUC极差（0.0063）；（3）在不平衡欺诈检测场景中，AUPRC不仅是精度指标，更是模型区分能力的必要条件——仅报告ROC-AUC将掩盖近一半的实际性能差异。

\subsubsection{结果分析}

从以上定量与可视化结果可以进一步提炼以下深层发现：

\begin{enumerate}
    \item \textbf{蒸馏改变深度表格模型的性能天花板}：DeepFM+KD（AUPRC=0.9495）以超过XGBoost教师1.64pp的幅度确立全局最优，验证了知识蒸馏的核心价值——将树模型的决策边界知识迁移至深度模型后，深度模型的额外参数容量可转化为真正的性能优势。FT-Transformer的基线AUPRC仅为0.9179，但FTT+KD跃升至0.9347（+1.68pp），与DeepFM+KD的增益（+1.59pp）高度一致——软标签蒸馏对注意力型和交互型两种异构架构表现出同等的正向迁移效应。
    \item \textbf{基座模型的归纳偏置对比}：在无蒸馏条件下，DeepFM（0.9336）$\approx$ XGBoost（0.9331）$>$ FT-Transformer（0.9179）。DeepFM的显式二阶交互在PCA空间中保有稳健性，而FT-Transformer的全局自注意力在缺乏跨特征语义关联的环境中出现性能退化——即便Sparsemax已将AUPRC从0.8738大幅回升至0.9179（+4.4pp），注意力分布仍近乎均匀（标准差0.0004），说明PCA匿名化对自注意力的影响无法仅通过归一化函数完全消除。
    \item \textbf{TreeReg的因果角色：方向明确但精确机制待确证}：FTT+KD（0.9347）$\rightarrow$ FTT+KD+TreeReg（0.8974）的-3.73pp断崖式下降，与"TreeReg中树先验对嵌入层的刚性约束与Sparsemax注意力的自由交互需求存在结构性冲突"的假设一致。但替代解释——$\lambda{=}0.1$对d\_model{=}48的嵌入空间过强、KD变体与基线的架构参数不完全一致（d\_model 48 vs 64, n\_layers 2 vs 3）——尚未通过控制实验排除。当前证据支持KD的安全有效性和TreeReg的警示信号，但对"结构性拮抗"的具体机制保持因果限定。
    \item \textbf{ROC-AUC的误导性}：全部模型ROC-AUC差异极小（0.9925--0.9988，极差仅0.0063），而AUPRC差异（0.8974--0.9495，极差0.0521）约为ROC-AUC极差的8.3倍。若仅报告ROC-AUC，DeepFM+KD比FTT+KD+TreeReg仅"好"0.0043，而AUPRC揭示的实际差距为0.0521。这进一步强化了在不平衡分类问题中使用AUPRC作为主要评估指标的必要性。
\end{enumerate}

\subsection{树监督蒸馏实验}

\subsubsection{方法设计}

第\ref{sec:experiments}节的结果表明，FT-Transformer的全局自注意力在引入Sparsemax后性能已大幅改善（AUPRC=0.9179），但其注意力权重仍近乎均匀——Sparsemax缓解了Softmax的完全塌缩（AUPRC提升4个百分点），但未能根本解决"哪些特征更重要"这一核心问题。XGBoost的轴对齐分裂天然知道特征的重要性排序（通过split gain），而FT-Transformer必须从零开始学习这一信息，在PCA匿名化特征空间中这一学习尤为困难。因此，本文的树监督蒸馏方案旨在弥补这一归纳偏置差距。

受Tree-Regularized Tabular Embeddings\cite{li2024tree}和Knowledge Distillation\cite{hinton2015distilling}等工作的启发，本文提出Tree-Supervised知识蒸馏方案，将XGBoost的树归纳偏置迁移至FT-Transformer。方案包含两个互补组件：

\textbf{方案3A：Leaf Soft Label Distillation（软标签蒸馏）。}将XGBoost的预测概率作为"软标签"，FT-Transformer以KL散度拟合教师的输出分布。总损失函数为：

\begin{equation}
\mathcal{L} = \alpha \cdot \text{FL}(y_s, y) + (1-\alpha) \cdot T^2 \cdot \text{KL}(p_s^T \parallel p_t^T)
\end{equation}

其中FL为Focal Loss（$\alpha=0.25, \gamma=2.0$），$T$为温度参数（$T=4.0$），$\alpha$为硬标签权重（$\alpha=0.7$）。温度软化使教师预测中的"暗知识"（如两个高概率类之间的微妙差异）得以传递，帮助学生学到比硬标签更丰富的决策边界信息。

\textbf{方案3B：Feature Split Importance Regularization（特征重要性正则化）。}提取XGBoost各特征的分裂增益（split gain），归一化为特征权重向量$w \in [0,1]^{29}$，将其作为FT-Transformer嵌入层的先验。正则化项为：

\begin{equation}
\mathcal{L}_{\text{reg}} = \lambda \sum_{i=1}^{29} (1 - w_i) \cdot \|\mathbf{E}_i\|_2
\end{equation}

其中$\mathbf{E}_i$为第$i$个特征的嵌入向量。该正则化鼓励模型为"低树重要性"特征分配较小嵌入范数，为"高树重要性"特征保留较大表达空间，从而将树模型的轴对齐特征选择偏好注入神经网络。

XGBoost特征权重（split gain归一化）显示，V14以$w_{14}=1.00$占据绝对主导，V12（$w_{12}=0.07$）和V4（$w_4=0.059$）为远距第二梯队，其余特征权重均低于0.02——与SHAP分析中V14的单特征主导现象一致。

\subsubsection{蒸馏实验结果}

回顾表\ref{tab:model_comparison}中蒸馏变体的定量数据，新增的FTT+KD（仅3A、不含TreeReg）消融实验彻底改变了蒸馏效应的因果归因：

\textbf{（1）知识蒸馏对两种架构均有显著增益。}FTT+KD以AUPRC=0.9347相较于FTT基线（0.9179）提升1.68个百分点，增益幅度与DeepFM+KD（+1.59pp）高度一致。这一平行结果推翻了此前"蒸馏是架构依赖的"的初步判断——软标签蒸馏对注意力型（FTT）和交互型（DeepFM）架构具有同等的正向迁移效应。此外，FTT+KD的校准近乎完美（ECE=0.0007，仅次于XGBoost教师），业务成本降至97（全局最低），且F1-Score=0.9279为所有模型最优。

\textbf{（2）TreeReg是退步的唯一元凶。}FTT+KD+TreeReg（3A+3B）的AUPRC仅为0.8974，不仅低于KD消融对照（0.9347），甚至低于未蒸馏的FTT基线（0.9179）。从FTT+KD（0.9347）到FTT+KD+TreeReg（0.8974）的-3.73pp断崖式下降，将因果链条唯一锁定于TreeReg（方案3B）：特征重要性正则化强制FT-Transformer的嵌入层模仿XGBoost的轴对齐特征偏好（V14权重=1.0，其余特征权重$<$0.07），这一硬性约束与Sparsemax注意力机制对自由特征交互学习的需求形成\textbf{结构性拮抗}——注意力需要基于PCA空间中特征间的实际语义关联分配权重，而TreeReg将其"锁定"在树模型的单一特征重要性排序上，导致嵌入空间的信息容量被急剧压缩。需要强调的是，3A+3B组合的良好ECE（0.0010）进一步印证了这一因果分析：软标签（3A）成功传递了决策边界平滑信息（ECE改善），但TreeReg（3B）对嵌入层的刚性约束抵消了KD带来的全部AUPRC增益。

\textbf{（3）校准质量的跨架构一致提升。}所有KD变体——FTT+KD（ECE=0.0007）、FTT+KD+TreeReg（ECE=0.0010）、DeepFM+KD（ECE=0.0008）——均取得卓越校准（ECE$<$0.002），验证了Hinton等人\cite{hinton2015distilling}的核心论点：软标签通过平滑类别边界，有效降低了深度模型的过度自信。这一效应与架构选择和TreeReg的存在均无关，是KD最稳健的收益。

\subsubsection{模型参数与归纳偏置}

各模型的实际参数量如下：FT-Transformer基线为153,921（KD变体因d\_model=48约为68K），DeepFM为161,648。XGBoost（300棵树$\times$深度8）的等效容量远低于深度模型，其基线AUPRC（0.9331）接近DeepFM（0.9336）但低于蒸馏变体（DeepFM+KD=0.9495, FTT+KD=0.9347）——验证了Grinsztajn等人\cite{grinsztajn2022tree}的观点：参数数量并非表格数据上性能的决定性因素，模型的归纳偏置才是关键。本文的蒸馏结果进一步表明，当深度模型通过知识蒸馏获得树模型的归纳偏置后，其额外的参数容量能够转化为真正的性能优势（DeepFM+KD超越XGBoost教师+1.64pp）。
